% ============================================================
% ACTIVIDAD PRÁCTICA: SIMULACIÓN DE PASE DE TURNO (MRM)
% Se puede usar sola o con % ============================================================
% ACTIVIDAD PRÁCTICA: SIMULACIÓN DE PASE DE TURNO (MRM)
% Se puede usar sola o con % ============================================================
% ACTIVIDAD PRÁCTICA: SIMULACIÓN DE PASE DE TURNO (MRM)
% Se puede usar sola o con % ============================================================
% ACTIVIDAD PRÁCTICA: SIMULACIÓN DE PASE DE TURNO (MRM)
% Se puede usar sola o con \input{actividad_mrm} dentro del manual.
% ============================================================
\clearpage
\section*{ACTIVIDAD PRÁCTICA: SIMULACIÓN DE PASE DE TURNO}
\subsection*{Factores humanos en mantenimiento (MRM)}

\begin{recuadro}
\textbf{Objetivo:} experimentar cómo una mala comunicación en el pase de turno genera errores
de mantenimiento, identificar los factores humanos que intervienen (la ``docena sucia'') y
aplicar una herramienta concreta del MRM, la planilla de pase de turno, para prevenirlos.
\end{recuadro}

\paragraph*{Duración:} 2 módulos (80 a 120 minutos).
\paragraph*{Lugar:} taller o laboratorio de aviónica.
\paragraph*{Materiales:}
\begin{itemize}
  \item Un equipo o panel por grupo, con tapas atornilladas y conectores desmontables
        (por ejemplo, una caja de instrumentos de práctica, una fuente o un equipo en desuso).
  \item Herramientas de mano y un recipiente para tornillería.
  \item Cinta o tarjetas de ``trabajo incompleto'' para señalizar.
  \item Copias de la \textbf{Planilla de pase de turno} y de la \textbf{Planilla de registro}
        (al final de esta actividad).
  \item Cronómetro.
\end{itemize}

% ------------------------------------------------------------
\subsubsection*{1. Introducción (15 minutos)}
El docente presenta brevemente los dos casos reales y la docena sucia.

\paragraph*{Continental Express 2574 (Texas, EE.\,UU., 1991).}
Durante el turno noche se comenzó a cambiar los bordes de ataque del estabilizador
horizontal de un Embraer 120. Se retiraron los tornillos superiores de uno de los lados,
pero el trabajo quedó sin terminar. El aviso en el pase de turno fue verbal y no llegó a
quien continuaba la tarea. El avión salió con esa pieza sujeta solo por los tornillos
inferiores; se desprendió en vuelo y murieron los 14 ocupantes.

\paragraph*{British Airways 5390 (Reino Unido, 1990).}
De noche y con poco personal, un supervisor reemplazó el parabrisas del piloto de un
BAC 1-11 eligiendo los tornillos ``a ojo'', sin consultar el catálogo de partes. La mayoría
eran levemente más finos que los especificados. En vuelo, la presión de cabina arrancó el
parabrisas; el comandante quedó parcialmente fuera del avión y fue sostenido por la
tripulación hasta el aterrizaje.

\begin{center}
\small
\textbf{La ``docena sucia'' (\textit{Dirty Dozen}) de factores humanos en mantenimiento}\\[3pt]
\begin{tabular}{|p{0.29\linewidth}|p{0.29\linewidth}|p{0.29\linewidth}|}
\hline
1. Falta de comunicación & 5. Falta de trabajo en equipo & 9. Falta de asertividad \\ \hline
2. Complacencia & 6. Fatiga & 10. Estrés \\ \hline
3. Falta de conocimiento & 7. Falta de recursos & 11. Falta de conciencia situacional \\ \hline
4. Distracción & 8. Presión & 12. Normas (malas costumbres del grupo) \\ \hline
\end{tabular}
\end{center}

% ------------------------------------------------------------
\subsubsection*{2. Organización}
El curso se divide en grupos de 4 a 6 integrantes. Cada grupo trabaja sobre su propio
equipo y se divide en dos ``turnos'':
\begin{itemize}
  \item \textbf{Turno A (saliente):} comienza el trabajo y lo deja a medias.
  \item \textbf{Turno B (entrante):} recibe el trabajo y debe terminarlo.
  \item \textbf{Un observador} por grupo (puede rotar): no participa del trabajo, registra lo
        que ocurre en la Planilla de registro.
\end{itemize}
El docente entrega al Turno A una \textbf{orden de trabajo} escrita, que el Turno B
\textbf{no ve}. Ejemplo de orden de trabajo:
\begin{enumerate}
  \item Retirar la tapa superior del equipo.
  \item Desconectar los conectores J1 y J3.
  \item Retirar el módulo interno indicado por el docente.
  \item Volver a instalar el módulo, reconectar J1 y J3 y colocar la tapa con todos sus
        tornillos.
  \item Verificar el funcionamiento.
\end{enumerate}

% ------------------------------------------------------------
\subsubsection*{3. Ronda 1: pase de turno ``informal'' (20 minutos)}
\begin{enumerate}
  \item El Turno A realiza los pasos 1 a 3 y, antes de terminar, el docente anuncia el
        \textbf{fin del turno}. El docente puede agregar, sin avisar, alguna condición extra
        (por ejemplo, dejar un tornillo de la tapa colocado sin ajustar).
  \item El Turno A tiene \textbf{un minuto} para contarle al Turno B, \textbf{solo de palabra},
        en qué quedó el trabajo. No pueden escribir ni mostrar el equipo.
  \item El Turno B termina el trabajo, \textbf{sin ver la orden de trabajo} y sin hacer
        preguntas al Turno A.
  \item El observador registra todo en la Planilla de registro.
\end{enumerate}

% ------------------------------------------------------------
\subsubsection*{4. Ronda 2: pase de turno con planilla (25 minutos)}
Se repite el ejercicio sobre el mismo equipo (o se intercambian los equipos entre grupos), y
\textbf{se invierten los roles}: quienes fueron Turno B ahora son Turno A. Esta vez:
\begin{enumerate}
  \item El Turno A completa la \textbf{Planilla de pase de turno} antes de irse y
        \textbf{señaliza} en el equipo lo que quedó abierto o desconectado.
  \item El pase se hace \textbf{cara a cara, frente al equipo}, leyendo la planilla. El Turno B
        puede hacer preguntas y firma que recibió el trabajo.
  \item El Turno B termina la tarea usando la planilla y la orden de trabajo, y realiza una
        \textbf{verificación final} antes de cerrar.
  \item El observador vuelve a completar la Planilla de registro.
\end{enumerate}

% ------------------------------------------------------------
\subsubsection*{5. Puesta en común (20 minutos)}
Cada grupo compara los resultados de las dos rondas y responde:
\begin{enumerate}
  \item ¿Cuántos errores u omisiones hubo en cada ronda? (tornillos faltantes o flojos,
        conectores sin reconectar o cruzados, pasos omitidos, verificación no realizada).
  \item ¿Qué factores de la docena sucia aparecieron en la Ronda 1? Den un ejemplo concreto de
        cada uno.
  \item ¿Qué cambió en la Ronda 2? ¿Qué parte de la planilla fue la más útil?
  \item ¿Alguien del Turno B notó que algo no cerraba y no lo dijo? ¿Por qué? (asertividad,
        presión de tiempo).
  \item ¿Qué tienen en común sus errores con los casos de Continental Express 2574 y
        British Airways 5390?
  \item Propongan una mejora para la planilla de pase de turno del taller de la escuela.
\end{enumerate}

% ------------------------------------------------------------
\subsubsection*{6. Evaluación}
\begin{itemize}
  \item Participación y cumplimiento del rol asignado en ambas rondas.
  \item Calidad y claridad de la Planilla de pase de turno completada.
  \item Correcta identificación de los factores humanos en la puesta en común.
  \item Respuestas escritas a las preguntas del punto 5 (entrega individual o grupal).
\end{itemize}

\begin{recuadro}
\textbf{Idea central:} en mantenimiento, \textbf{lo que no está escrito no está hecho}. El
pase de turno se hace por escrito, cara a cara y frente al trabajo; ante la duda se verifica
con la documentación aprobada y nunca se supone.
\end{recuadro}

% ============================================================
% PLANILLA DE PASE DE TURNO (para imprimir)
% ============================================================
\clearpage
\begin{center}
{\large\bfseries PLANILLA DE PASE DE TURNO}\\[2pt]
{\small Taller de Aviónica -- E.E.S.T. N° 7 ``Taller Regional Quilmes''}
\end{center}

\renewcommand{\arraystretch}{1.4}
\noindent
\begin{tabular}{|>{\raggedright\arraybackslash}p{0.34\linewidth}|p{0.595\linewidth}|}
\hline
\textbf{Fecha} & \\ \hline
\textbf{Equipo / aeronave} & \\ \hline
\textbf{N° de orden de trabajo} & \\ \hline
\textbf{Turno saliente (nombres)} & \\ \hline
\textbf{Turno entrante (nombres)} & \\ \hline
\end{tabular}

\vspace{5pt}
\noindent\textbf{1. Tareas realizadas} (con el paso de la orden de trabajo)\\[2pt]
\begin{tabular}{|p{0.10\linewidth}|p{0.60\linewidth}|p{0.215\linewidth}|}
\hline
\textbf{Paso} & \textbf{Descripción de lo realizado} & \textbf{Firma} \\ \hline
 & & \\ \hline
 & & \\ \hline
 & & \\ \hline
\end{tabular}

\vspace{5pt}
\noindent\textbf{2. Tareas pendientes}\\[2pt]
\begin{tabular}{|p{0.10\linewidth}|p{0.83\linewidth}|}
\hline
\textbf{Paso} & \textbf{Qué falta hacer} \\ \hline
 & \\ \hline
 & \\ \hline
 & \\ \hline
\end{tabular}

\vspace{5pt}
\noindent\textbf{3. Estado en que queda el equipo}\\[2pt]
\begin{tabular}{|p{0.45\linewidth}|p{0.48\linewidth}|}
\hline
\textbf{Desmontado / desconectado} & \textbf{Dónde quedó (ubicación, rótulo)} \\ \hline
 & \\ \hline
 & \\ \hline
\textbf{Tornillería / fijaciones retiradas} & \textbf{Cantidad y dónde están} \\ \hline
 & \\ \hline
\end{tabular}

\vspace{5pt}
\noindent
\begin{tabular}{|p{0.62\linewidth}|c|c|}
\hline
\textbf{4. Verificaciones} & \textbf{Sí} & \textbf{No} \\ \hline
¿Se señalizó en el equipo que el trabajo está incompleto? & \hspace{0.6cm} & \hspace{0.6cm} \\ \hline
¿Hay herramientas o elementos sueltos dentro del equipo? & & \\ \hline
¿Se informó alguna anomalía o duda? (detallar abajo) & & \\ \hline
\end{tabular}

\vspace{5pt}
\noindent\textbf{5. Observaciones, anomalías o advertencias}\\[2pt]
\begin{tabular}{|p{0.955\linewidth}|}
\hline
\\ \hline
\\ \hline
\end{tabular}

\vspace{10pt}
\noindent
\begin{tabular}{|p{0.46\linewidth}|p{0.46\linewidth}|}
\hline
\textbf{Entrega (turno saliente)} & \textbf{Recibe (turno entrante)} \\
Firma y aclaración: & Firma y aclaración: \\[14pt] \hline
\end{tabular}
\renewcommand{\arraystretch}{1.15}

% ============================================================
% PLANILLA DE REGISTRO DEL OBSERVADOR (para imprimir)
% ============================================================
\clearpage
\begin{center}
{\large\bfseries PLANILLA DE REGISTRO DEL OBSERVADOR}\\[2pt]
{\small Simulación de pase de turno -- MRM}
\end{center}

\renewcommand{\arraystretch}{1.6}
\noindent
\begin{tabular}{|>{\raggedright\arraybackslash}p{0.34\linewidth}|p{0.595\linewidth}|}
\hline
\textbf{Grupo} & \\ \hline
\textbf{Observador/a} & \\ \hline
\end{tabular}

\vspace{8pt}
\noindent\textbf{Errores y omisiones detectados}\\[2pt]
\begin{tabular}{|p{0.50\linewidth}|>{\centering\arraybackslash}p{0.19\linewidth}|>{\centering\arraybackslash}p{0.19\linewidth}|}
\hline
\textbf{Tipo de error} & \textbf{Ronda 1} & \textbf{Ronda 2} \\ \hline
Tornillos faltantes, flojos o sobrantes & & \\ \hline
Conectores sin reconectar o cruzados & & \\ \hline
Pasos de la orden de trabajo omitidos & & \\ \hline
Componente mal instalado & & \\ \hline
Herramienta o elemento olvidado en el equipo & & \\ \hline
Verificación final no realizada & & \\ \hline
\textbf{Total} & & \\ \hline
\end{tabular}

\vspace{8pt}
\noindent\textbf{Factores de la docena sucia observados} (marcar y anotar un ejemplo)\\[2pt]
\begin{tabular}{|p{0.36\linewidth}|>{\centering\arraybackslash}p{0.08\linewidth}|>{\centering\arraybackslash}p{0.08\linewidth}|p{0.34\linewidth}|}
\hline
\textbf{Factor} & \textbf{R1} & \textbf{R2} & \textbf{Ejemplo observado} \\ \hline
1. Falta de comunicación & & & \\ \hline
2. Complacencia & & & \\ \hline
3. Falta de conocimiento & & & \\ \hline
4. Distracción & & & \\ \hline
5. Falta de trabajo en equipo & & & \\ \hline
6. Fatiga & & & \\ \hline
7. Falta de recursos & & & \\ \hline
8. Presión & & & \\ \hline
9. Falta de asertividad & & & \\ \hline
10. Estrés & & & \\ \hline
11. Falta de conciencia situacional & & & \\ \hline
12. Normas & & & \\ \hline
\end{tabular}

\vspace{8pt}
\noindent\textbf{Comentarios del observador}\\[2pt]
\begin{tabular}{|p{0.955\linewidth}|}
\hline
\\ \hline
\\ \hline
\\ \hline
\end{tabular}
\renewcommand{\arraystretch}{1.15}
 dentro del manual.
% ============================================================
\clearpage
\section*{ACTIVIDAD PRÁCTICA: SIMULACIÓN DE PASE DE TURNO}
\subsection*{Factores humanos en mantenimiento (MRM)}

\begin{recuadro}
\textbf{Objetivo:} experimentar cómo una mala comunicación en el pase de turno genera errores
de mantenimiento, identificar los factores humanos que intervienen (la ``docena sucia'') y
aplicar una herramienta concreta del MRM, la planilla de pase de turno, para prevenirlos.
\end{recuadro}

\paragraph*{Duración:} 2 módulos (80 a 120 minutos).
\paragraph*{Lugar:} taller o laboratorio de aviónica.
\paragraph*{Materiales:}
\begin{itemize}
  \item Un equipo o panel por grupo, con tapas atornilladas y conectores desmontables
        (por ejemplo, una caja de instrumentos de práctica, una fuente o un equipo en desuso).
  \item Herramientas de mano y un recipiente para tornillería.
  \item Cinta o tarjetas de ``trabajo incompleto'' para señalizar.
  \item Copias de la \textbf{Planilla de pase de turno} y de la \textbf{Planilla de registro}
        (al final de esta actividad).
  \item Cronómetro.
\end{itemize}

% ------------------------------------------------------------
\subsubsection*{1. Introducción (15 minutos)}
El docente presenta brevemente los dos casos reales y la docena sucia.

\paragraph*{Continental Express 2574 (Texas, EE.\,UU., 1991).}
Durante el turno noche se comenzó a cambiar los bordes de ataque del estabilizador
horizontal de un Embraer 120. Se retiraron los tornillos superiores de uno de los lados,
pero el trabajo quedó sin terminar. El aviso en el pase de turno fue verbal y no llegó a
quien continuaba la tarea. El avión salió con esa pieza sujeta solo por los tornillos
inferiores; se desprendió en vuelo y murieron los 14 ocupantes.

\paragraph*{British Airways 5390 (Reino Unido, 1990).}
De noche y con poco personal, un supervisor reemplazó el parabrisas del piloto de un
BAC 1-11 eligiendo los tornillos ``a ojo'', sin consultar el catálogo de partes. La mayoría
eran levemente más finos que los especificados. En vuelo, la presión de cabina arrancó el
parabrisas; el comandante quedó parcialmente fuera del avión y fue sostenido por la
tripulación hasta el aterrizaje.

\begin{center}
\small
\textbf{La ``docena sucia'' (\textit{Dirty Dozen}) de factores humanos en mantenimiento}\\[3pt]
\begin{tabular}{|p{0.29\linewidth}|p{0.29\linewidth}|p{0.29\linewidth}|}
\hline
1. Falta de comunicación & 5. Falta de trabajo en equipo & 9. Falta de asertividad \\ \hline
2. Complacencia & 6. Fatiga & 10. Estrés \\ \hline
3. Falta de conocimiento & 7. Falta de recursos & 11. Falta de conciencia situacional \\ \hline
4. Distracción & 8. Presión & 12. Normas (malas costumbres del grupo) \\ \hline
\end{tabular}
\end{center}

% ------------------------------------------------------------
\subsubsection*{2. Organización}
El curso se divide en grupos de 4 a 6 integrantes. Cada grupo trabaja sobre su propio
equipo y se divide en dos ``turnos'':
\begin{itemize}
  \item \textbf{Turno A (saliente):} comienza el trabajo y lo deja a medias.
  \item \textbf{Turno B (entrante):} recibe el trabajo y debe terminarlo.
  \item \textbf{Un observador} por grupo (puede rotar): no participa del trabajo, registra lo
        que ocurre en la Planilla de registro.
\end{itemize}
El docente entrega al Turno A una \textbf{orden de trabajo} escrita, que el Turno B
\textbf{no ve}. Ejemplo de orden de trabajo:
\begin{enumerate}
  \item Retirar la tapa superior del equipo.
  \item Desconectar los conectores J1 y J3.
  \item Retirar el módulo interno indicado por el docente.
  \item Volver a instalar el módulo, reconectar J1 y J3 y colocar la tapa con todos sus
        tornillos.
  \item Verificar el funcionamiento.
\end{enumerate}

% ------------------------------------------------------------
\subsubsection*{3. Ronda 1: pase de turno ``informal'' (20 minutos)}
\begin{enumerate}
  \item El Turno A realiza los pasos 1 a 3 y, antes de terminar, el docente anuncia el
        \textbf{fin del turno}. El docente puede agregar, sin avisar, alguna condición extra
        (por ejemplo, dejar un tornillo de la tapa colocado sin ajustar).
  \item El Turno A tiene \textbf{un minuto} para contarle al Turno B, \textbf{solo de palabra},
        en qué quedó el trabajo. No pueden escribir ni mostrar el equipo.
  \item El Turno B termina el trabajo, \textbf{sin ver la orden de trabajo} y sin hacer
        preguntas al Turno A.
  \item El observador registra todo en la Planilla de registro.
\end{enumerate}

% ------------------------------------------------------------
\subsubsection*{4. Ronda 2: pase de turno con planilla (25 minutos)}
Se repite el ejercicio sobre el mismo equipo (o se intercambian los equipos entre grupos), y
\textbf{se invierten los roles}: quienes fueron Turno B ahora son Turno A. Esta vez:
\begin{enumerate}
  \item El Turno A completa la \textbf{Planilla de pase de turno} antes de irse y
        \textbf{señaliza} en el equipo lo que quedó abierto o desconectado.
  \item El pase se hace \textbf{cara a cara, frente al equipo}, leyendo la planilla. El Turno B
        puede hacer preguntas y firma que recibió el trabajo.
  \item El Turno B termina la tarea usando la planilla y la orden de trabajo, y realiza una
        \textbf{verificación final} antes de cerrar.
  \item El observador vuelve a completar la Planilla de registro.
\end{enumerate}

% ------------------------------------------------------------
\subsubsection*{5. Puesta en común (20 minutos)}
Cada grupo compara los resultados de las dos rondas y responde:
\begin{enumerate}
  \item ¿Cuántos errores u omisiones hubo en cada ronda? (tornillos faltantes o flojos,
        conectores sin reconectar o cruzados, pasos omitidos, verificación no realizada).
  \item ¿Qué factores de la docena sucia aparecieron en la Ronda 1? Den un ejemplo concreto de
        cada uno.
  \item ¿Qué cambió en la Ronda 2? ¿Qué parte de la planilla fue la más útil?
  \item ¿Alguien del Turno B notó que algo no cerraba y no lo dijo? ¿Por qué? (asertividad,
        presión de tiempo).
  \item ¿Qué tienen en común sus errores con los casos de Continental Express 2574 y
        British Airways 5390?
  \item Propongan una mejora para la planilla de pase de turno del taller de la escuela.
\end{enumerate}

% ------------------------------------------------------------
\subsubsection*{6. Evaluación}
\begin{itemize}
  \item Participación y cumplimiento del rol asignado en ambas rondas.
  \item Calidad y claridad de la Planilla de pase de turno completada.
  \item Correcta identificación de los factores humanos en la puesta en común.
  \item Respuestas escritas a las preguntas del punto 5 (entrega individual o grupal).
\end{itemize}

\begin{recuadro}
\textbf{Idea central:} en mantenimiento, \textbf{lo que no está escrito no está hecho}. El
pase de turno se hace por escrito, cara a cara y frente al trabajo; ante la duda se verifica
con la documentación aprobada y nunca se supone.
\end{recuadro}

% ============================================================
% PLANILLA DE PASE DE TURNO (para imprimir)
% ============================================================
\clearpage
\begin{center}
{\large\bfseries PLANILLA DE PASE DE TURNO}\\[2pt]
{\small Taller de Aviónica -- E.E.S.T. N° 7 ``Taller Regional Quilmes''}
\end{center}

\renewcommand{\arraystretch}{1.4}
\noindent
\begin{tabular}{|>{\raggedright\arraybackslash}p{0.34\linewidth}|p{0.595\linewidth}|}
\hline
\textbf{Fecha} & \\ \hline
\textbf{Equipo / aeronave} & \\ \hline
\textbf{N° de orden de trabajo} & \\ \hline
\textbf{Turno saliente (nombres)} & \\ \hline
\textbf{Turno entrante (nombres)} & \\ \hline
\end{tabular}

\vspace{5pt}
\noindent\textbf{1. Tareas realizadas} (con el paso de la orden de trabajo)\\[2pt]
\begin{tabular}{|p{0.10\linewidth}|p{0.60\linewidth}|p{0.215\linewidth}|}
\hline
\textbf{Paso} & \textbf{Descripción de lo realizado} & \textbf{Firma} \\ \hline
 & & \\ \hline
 & & \\ \hline
 & & \\ \hline
\end{tabular}

\vspace{5pt}
\noindent\textbf{2. Tareas pendientes}\\[2pt]
\begin{tabular}{|p{0.10\linewidth}|p{0.83\linewidth}|}
\hline
\textbf{Paso} & \textbf{Qué falta hacer} \\ \hline
 & \\ \hline
 & \\ \hline
 & \\ \hline
\end{tabular}

\vspace{5pt}
\noindent\textbf{3. Estado en que queda el equipo}\\[2pt]
\begin{tabular}{|p{0.45\linewidth}|p{0.48\linewidth}|}
\hline
\textbf{Desmontado / desconectado} & \textbf{Dónde quedó (ubicación, rótulo)} \\ \hline
 & \\ \hline
 & \\ \hline
\textbf{Tornillería / fijaciones retiradas} & \textbf{Cantidad y dónde están} \\ \hline
 & \\ \hline
\end{tabular}

\vspace{5pt}
\noindent
\begin{tabular}{|p{0.62\linewidth}|c|c|}
\hline
\textbf{4. Verificaciones} & \textbf{Sí} & \textbf{No} \\ \hline
¿Se señalizó en el equipo que el trabajo está incompleto? & \hspace{0.6cm} & \hspace{0.6cm} \\ \hline
¿Hay herramientas o elementos sueltos dentro del equipo? & & \\ \hline
¿Se informó alguna anomalía o duda? (detallar abajo) & & \\ \hline
\end{tabular}

\vspace{5pt}
\noindent\textbf{5. Observaciones, anomalías o advertencias}\\[2pt]
\begin{tabular}{|p{0.955\linewidth}|}
\hline
\\ \hline
\\ \hline
\end{tabular}

\vspace{10pt}
\noindent
\begin{tabular}{|p{0.46\linewidth}|p{0.46\linewidth}|}
\hline
\textbf{Entrega (turno saliente)} & \textbf{Recibe (turno entrante)} \\
Firma y aclaración: & Firma y aclaración: \\[14pt] \hline
\end{tabular}
\renewcommand{\arraystretch}{1.15}

% ============================================================
% PLANILLA DE REGISTRO DEL OBSERVADOR (para imprimir)
% ============================================================
\clearpage
\begin{center}
{\large\bfseries PLANILLA DE REGISTRO DEL OBSERVADOR}\\[2pt]
{\small Simulación de pase de turno -- MRM}
\end{center}

\renewcommand{\arraystretch}{1.6}
\noindent
\begin{tabular}{|>{\raggedright\arraybackslash}p{0.34\linewidth}|p{0.595\linewidth}|}
\hline
\textbf{Grupo} & \\ \hline
\textbf{Observador/a} & \\ \hline
\end{tabular}

\vspace{8pt}
\noindent\textbf{Errores y omisiones detectados}\\[2pt]
\begin{tabular}{|p{0.50\linewidth}|>{\centering\arraybackslash}p{0.19\linewidth}|>{\centering\arraybackslash}p{0.19\linewidth}|}
\hline
\textbf{Tipo de error} & \textbf{Ronda 1} & \textbf{Ronda 2} \\ \hline
Tornillos faltantes, flojos o sobrantes & & \\ \hline
Conectores sin reconectar o cruzados & & \\ \hline
Pasos de la orden de trabajo omitidos & & \\ \hline
Componente mal instalado & & \\ \hline
Herramienta o elemento olvidado en el equipo & & \\ \hline
Verificación final no realizada & & \\ \hline
\textbf{Total} & & \\ \hline
\end{tabular}

\vspace{8pt}
\noindent\textbf{Factores de la docena sucia observados} (marcar y anotar un ejemplo)\\[2pt]
\begin{tabular}{|p{0.36\linewidth}|>{\centering\arraybackslash}p{0.08\linewidth}|>{\centering\arraybackslash}p{0.08\linewidth}|p{0.34\linewidth}|}
\hline
\textbf{Factor} & \textbf{R1} & \textbf{R2} & \textbf{Ejemplo observado} \\ \hline
1. Falta de comunicación & & & \\ \hline
2. Complacencia & & & \\ \hline
3. Falta de conocimiento & & & \\ \hline
4. Distracción & & & \\ \hline
5. Falta de trabajo en equipo & & & \\ \hline
6. Fatiga & & & \\ \hline
7. Falta de recursos & & & \\ \hline
8. Presión & & & \\ \hline
9. Falta de asertividad & & & \\ \hline
10. Estrés & & & \\ \hline
11. Falta de conciencia situacional & & & \\ \hline
12. Normas & & & \\ \hline
\end{tabular}

\vspace{8pt}
\noindent\textbf{Comentarios del observador}\\[2pt]
\begin{tabular}{|p{0.955\linewidth}|}
\hline
\\ \hline
\\ \hline
\\ \hline
\end{tabular}
\renewcommand{\arraystretch}{1.15}
 dentro del manual.
% ============================================================
\clearpage
\section*{ACTIVIDAD PRÁCTICA: SIMULACIÓN DE PASE DE TURNO}
\subsection*{Factores humanos en mantenimiento (MRM)}

\begin{recuadro}
\textbf{Objetivo:} experimentar cómo una mala comunicación en el pase de turno genera errores
de mantenimiento, identificar los factores humanos que intervienen (la ``docena sucia'') y
aplicar una herramienta concreta del MRM, la planilla de pase de turno, para prevenirlos.
\end{recuadro}

\paragraph*{Duración:} 2 módulos (80 a 120 minutos).
\paragraph*{Lugar:} taller o laboratorio de aviónica.
\paragraph*{Materiales:}
\begin{itemize}
  \item Un equipo o panel por grupo, con tapas atornilladas y conectores desmontables
        (por ejemplo, una caja de instrumentos de práctica, una fuente o un equipo en desuso).
  \item Herramientas de mano y un recipiente para tornillería.
  \item Cinta o tarjetas de ``trabajo incompleto'' para señalizar.
  \item Copias de la \textbf{Planilla de pase de turno} y de la \textbf{Planilla de registro}
        (al final de esta actividad).
  \item Cronómetro.
\end{itemize}

% ------------------------------------------------------------
\subsubsection*{1. Introducción (15 minutos)}
El docente presenta brevemente los dos casos reales y la docena sucia.

\paragraph*{Continental Express 2574 (Texas, EE.\,UU., 1991).}
Durante el turno noche se comenzó a cambiar los bordes de ataque del estabilizador
horizontal de un Embraer 120. Se retiraron los tornillos superiores de uno de los lados,
pero el trabajo quedó sin terminar. El aviso en el pase de turno fue verbal y no llegó a
quien continuaba la tarea. El avión salió con esa pieza sujeta solo por los tornillos
inferiores; se desprendió en vuelo y murieron los 14 ocupantes.

\paragraph*{British Airways 5390 (Reino Unido, 1990).}
De noche y con poco personal, un supervisor reemplazó el parabrisas del piloto de un
BAC 1-11 eligiendo los tornillos ``a ojo'', sin consultar el catálogo de partes. La mayoría
eran levemente más finos que los especificados. En vuelo, la presión de cabina arrancó el
parabrisas; el comandante quedó parcialmente fuera del avión y fue sostenido por la
tripulación hasta el aterrizaje.

\begin{center}
\small
\textbf{La ``docena sucia'' (\textit{Dirty Dozen}) de factores humanos en mantenimiento}\\[3pt]
\begin{tabular}{|p{0.29\linewidth}|p{0.29\linewidth}|p{0.29\linewidth}|}
\hline
1. Falta de comunicación & 5. Falta de trabajo en equipo & 9. Falta de asertividad \\ \hline
2. Complacencia & 6. Fatiga & 10. Estrés \\ \hline
3. Falta de conocimiento & 7. Falta de recursos & 11. Falta de conciencia situacional \\ \hline
4. Distracción & 8. Presión & 12. Normas (malas costumbres del grupo) \\ \hline
\end{tabular}
\end{center}

% ------------------------------------------------------------
\subsubsection*{2. Organización}
El curso se divide en grupos de 4 a 6 integrantes. Cada grupo trabaja sobre su propio
equipo y se divide en dos ``turnos'':
\begin{itemize}
  \item \textbf{Turno A (saliente):} comienza el trabajo y lo deja a medias.
  \item \textbf{Turno B (entrante):} recibe el trabajo y debe terminarlo.
  \item \textbf{Un observador} por grupo (puede rotar): no participa del trabajo, registra lo
        que ocurre en la Planilla de registro.
\end{itemize}
El docente entrega al Turno A una \textbf{orden de trabajo} escrita, que el Turno B
\textbf{no ve}. Ejemplo de orden de trabajo:
\begin{enumerate}
  \item Retirar la tapa superior del equipo.
  \item Desconectar los conectores J1 y J3.
  \item Retirar el módulo interno indicado por el docente.
  \item Volver a instalar el módulo, reconectar J1 y J3 y colocar la tapa con todos sus
        tornillos.
  \item Verificar el funcionamiento.
\end{enumerate}

% ------------------------------------------------------------
\subsubsection*{3. Ronda 1: pase de turno ``informal'' (20 minutos)}
\begin{enumerate}
  \item El Turno A realiza los pasos 1 a 3 y, antes de terminar, el docente anuncia el
        \textbf{fin del turno}. El docente puede agregar, sin avisar, alguna condición extra
        (por ejemplo, dejar un tornillo de la tapa colocado sin ajustar).
  \item El Turno A tiene \textbf{un minuto} para contarle al Turno B, \textbf{solo de palabra},
        en qué quedó el trabajo. No pueden escribir ni mostrar el equipo.
  \item El Turno B termina el trabajo, \textbf{sin ver la orden de trabajo} y sin hacer
        preguntas al Turno A.
  \item El observador registra todo en la Planilla de registro.
\end{enumerate}

% ------------------------------------------------------------
\subsubsection*{4. Ronda 2: pase de turno con planilla (25 minutos)}
Se repite el ejercicio sobre el mismo equipo (o se intercambian los equipos entre grupos), y
\textbf{se invierten los roles}: quienes fueron Turno B ahora son Turno A. Esta vez:
\begin{enumerate}
  \item El Turno A completa la \textbf{Planilla de pase de turno} antes de irse y
        \textbf{señaliza} en el equipo lo que quedó abierto o desconectado.
  \item El pase se hace \textbf{cara a cara, frente al equipo}, leyendo la planilla. El Turno B
        puede hacer preguntas y firma que recibió el trabajo.
  \item El Turno B termina la tarea usando la planilla y la orden de trabajo, y realiza una
        \textbf{verificación final} antes de cerrar.
  \item El observador vuelve a completar la Planilla de registro.
\end{enumerate}

% ------------------------------------------------------------
\subsubsection*{5. Puesta en común (20 minutos)}
Cada grupo compara los resultados de las dos rondas y responde:
\begin{enumerate}
  \item ¿Cuántos errores u omisiones hubo en cada ronda? (tornillos faltantes o flojos,
        conectores sin reconectar o cruzados, pasos omitidos, verificación no realizada).
  \item ¿Qué factores de la docena sucia aparecieron en la Ronda 1? Den un ejemplo concreto de
        cada uno.
  \item ¿Qué cambió en la Ronda 2? ¿Qué parte de la planilla fue la más útil?
  \item ¿Alguien del Turno B notó que algo no cerraba y no lo dijo? ¿Por qué? (asertividad,
        presión de tiempo).
  \item ¿Qué tienen en común sus errores con los casos de Continental Express 2574 y
        British Airways 5390?
  \item Propongan una mejora para la planilla de pase de turno del taller de la escuela.
\end{enumerate}

% ------------------------------------------------------------
\subsubsection*{6. Evaluación}
\begin{itemize}
  \item Participación y cumplimiento del rol asignado en ambas rondas.
  \item Calidad y claridad de la Planilla de pase de turno completada.
  \item Correcta identificación de los factores humanos en la puesta en común.
  \item Respuestas escritas a las preguntas del punto 5 (entrega individual o grupal).
\end{itemize}

\begin{recuadro}
\textbf{Idea central:} en mantenimiento, \textbf{lo que no está escrito no está hecho}. El
pase de turno se hace por escrito, cara a cara y frente al trabajo; ante la duda se verifica
con la documentación aprobada y nunca se supone.
\end{recuadro}

% ============================================================
% PLANILLA DE PASE DE TURNO (para imprimir)
% ============================================================
\clearpage
\begin{center}
{\large\bfseries PLANILLA DE PASE DE TURNO}\\[2pt]
{\small Taller de Aviónica -- E.E.S.T. N° 7 ``Taller Regional Quilmes''}
\end{center}

\renewcommand{\arraystretch}{1.4}
\noindent
\begin{tabular}{|>{\raggedright\arraybackslash}p{0.34\linewidth}|p{0.595\linewidth}|}
\hline
\textbf{Fecha} & \\ \hline
\textbf{Equipo / aeronave} & \\ \hline
\textbf{N° de orden de trabajo} & \\ \hline
\textbf{Turno saliente (nombres)} & \\ \hline
\textbf{Turno entrante (nombres)} & \\ \hline
\end{tabular}

\vspace{5pt}
\noindent\textbf{1. Tareas realizadas} (con el paso de la orden de trabajo)\\[2pt]
\begin{tabular}{|p{0.10\linewidth}|p{0.60\linewidth}|p{0.215\linewidth}|}
\hline
\textbf{Paso} & \textbf{Descripción de lo realizado} & \textbf{Firma} \\ \hline
 & & \\ \hline
 & & \\ \hline
 & & \\ \hline
\end{tabular}

\vspace{5pt}
\noindent\textbf{2. Tareas pendientes}\\[2pt]
\begin{tabular}{|p{0.10\linewidth}|p{0.83\linewidth}|}
\hline
\textbf{Paso} & \textbf{Qué falta hacer} \\ \hline
 & \\ \hline
 & \\ \hline
 & \\ \hline
\end{tabular}

\vspace{5pt}
\noindent\textbf{3. Estado en que queda el equipo}\\[2pt]
\begin{tabular}{|p{0.45\linewidth}|p{0.48\linewidth}|}
\hline
\textbf{Desmontado / desconectado} & \textbf{Dónde quedó (ubicación, rótulo)} \\ \hline
 & \\ \hline
 & \\ \hline
\textbf{Tornillería / fijaciones retiradas} & \textbf{Cantidad y dónde están} \\ \hline
 & \\ \hline
\end{tabular}

\vspace{5pt}
\noindent
\begin{tabular}{|p{0.62\linewidth}|c|c|}
\hline
\textbf{4. Verificaciones} & \textbf{Sí} & \textbf{No} \\ \hline
¿Se señalizó en el equipo que el trabajo está incompleto? & \hspace{0.6cm} & \hspace{0.6cm} \\ \hline
¿Hay herramientas o elementos sueltos dentro del equipo? & & \\ \hline
¿Se informó alguna anomalía o duda? (detallar abajo) & & \\ \hline
\end{tabular}

\vspace{5pt}
\noindent\textbf{5. Observaciones, anomalías o advertencias}\\[2pt]
\begin{tabular}{|p{0.955\linewidth}|}
\hline
\\ \hline
\\ \hline
\end{tabular}

\vspace{10pt}
\noindent
\begin{tabular}{|p{0.46\linewidth}|p{0.46\linewidth}|}
\hline
\textbf{Entrega (turno saliente)} & \textbf{Recibe (turno entrante)} \\
Firma y aclaración: & Firma y aclaración: \\[14pt] \hline
\end{tabular}
\renewcommand{\arraystretch}{1.15}

% ============================================================
% PLANILLA DE REGISTRO DEL OBSERVADOR (para imprimir)
% ============================================================
\clearpage
\begin{center}
{\large\bfseries PLANILLA DE REGISTRO DEL OBSERVADOR}\\[2pt]
{\small Simulación de pase de turno -- MRM}
\end{center}

\renewcommand{\arraystretch}{1.6}
\noindent
\begin{tabular}{|>{\raggedright\arraybackslash}p{0.34\linewidth}|p{0.595\linewidth}|}
\hline
\textbf{Grupo} & \\ \hline
\textbf{Observador/a} & \\ \hline
\end{tabular}

\vspace{8pt}
\noindent\textbf{Errores y omisiones detectados}\\[2pt]
\begin{tabular}{|p{0.50\linewidth}|>{\centering\arraybackslash}p{0.19\linewidth}|>{\centering\arraybackslash}p{0.19\linewidth}|}
\hline
\textbf{Tipo de error} & \textbf{Ronda 1} & \textbf{Ronda 2} \\ \hline
Tornillos faltantes, flojos o sobrantes & & \\ \hline
Conectores sin reconectar o cruzados & & \\ \hline
Pasos de la orden de trabajo omitidos & & \\ \hline
Componente mal instalado & & \\ \hline
Herramienta o elemento olvidado en el equipo & & \\ \hline
Verificación final no realizada & & \\ \hline
\textbf{Total} & & \\ \hline
\end{tabular}

\vspace{8pt}
\noindent\textbf{Factores de la docena sucia observados} (marcar y anotar un ejemplo)\\[2pt]
\begin{tabular}{|p{0.36\linewidth}|>{\centering\arraybackslash}p{0.08\linewidth}|>{\centering\arraybackslash}p{0.08\linewidth}|p{0.34\linewidth}|}
\hline
\textbf{Factor} & \textbf{R1} & \textbf{R2} & \textbf{Ejemplo observado} \\ \hline
1. Falta de comunicación & & & \\ \hline
2. Complacencia & & & \\ \hline
3. Falta de conocimiento & & & \\ \hline
4. Distracción & & & \\ \hline
5. Falta de trabajo en equipo & & & \\ \hline
6. Fatiga & & & \\ \hline
7. Falta de recursos & & & \\ \hline
8. Presión & & & \\ \hline
9. Falta de asertividad & & & \\ \hline
10. Estrés & & & \\ \hline
11. Falta de conciencia situacional & & & \\ \hline
12. Normas & & & \\ \hline
\end{tabular}

\vspace{8pt}
\noindent\textbf{Comentarios del observador}\\[2pt]
\begin{tabular}{|p{0.955\linewidth}|}
\hline
\\ \hline
\\ \hline
\\ \hline
\end{tabular}
\renewcommand{\arraystretch}{1.15}
 dentro del manual.
% ============================================================
\clearpage
\section*{ACTIVIDAD PRÁCTICA: SIMULACIÓN DE PASE DE TURNO}
\subsection*{Factores humanos en mantenimiento (MRM)}

\begin{recuadro}
\textbf{Objetivo:} experimentar cómo una mala comunicación en el pase de turno genera errores
de mantenimiento, identificar los factores humanos que intervienen (la ``docena sucia'') y
aplicar una herramienta concreta del MRM, la planilla de pase de turno, para prevenirlos.
\end{recuadro}

\paragraph*{Duración:} 2 módulos (80 a 120 minutos).
\paragraph*{Lugar:} taller o laboratorio de aviónica.
\paragraph*{Materiales:}
\begin{itemize}
  \item Un equipo o panel por grupo, con tapas atornilladas y conectores desmontables
        (por ejemplo, una caja de instrumentos de práctica, una fuente o un equipo en desuso).
  \item Herramientas de mano y un recipiente para tornillería.
  \item Cinta o tarjetas de ``trabajo incompleto'' para señalizar.
  \item Copias de la \textbf{Planilla de pase de turno} y de la \textbf{Planilla de registro}
        (al final de esta actividad).
  \item Cronómetro.
\end{itemize}

% ------------------------------------------------------------
\subsubsection*{1. Introducción (15 minutos)}
El docente presenta brevemente los dos casos reales y la docena sucia.

\paragraph*{Continental Express 2574 (Texas, EE.\,UU., 1991).}
Durante el turno noche se comenzó a cambiar los bordes de ataque del estabilizador
horizontal de un Embraer 120. Se retiraron los tornillos superiores de uno de los lados,
pero el trabajo quedó sin terminar. El aviso en el pase de turno fue verbal y no llegó a
quien continuaba la tarea. El avión salió con esa pieza sujeta solo por los tornillos
inferiores; se desprendió en vuelo y murieron los 14 ocupantes.

\paragraph*{British Airways 5390 (Reino Unido, 1990).}
De noche y con poco personal, un supervisor reemplazó el parabrisas del piloto de un
BAC 1-11 eligiendo los tornillos ``a ojo'', sin consultar el catálogo de partes. La mayoría
eran levemente más finos que los especificados. En vuelo, la presión de cabina arrancó el
parabrisas; el comandante quedó parcialmente fuera del avión y fue sostenido por la
tripulación hasta el aterrizaje.

\begin{center}
\small
\textbf{La ``docena sucia'' (\textit{Dirty Dozen}) de factores humanos en mantenimiento}\\[3pt]
\begin{tabular}{|p{0.29\linewidth}|p{0.29\linewidth}|p{0.29\linewidth}|}
\hline
1. Falta de comunicación & 5. Falta de trabajo en equipo & 9. Falta de asertividad \\ \hline
2. Complacencia & 6. Fatiga & 10. Estrés \\ \hline
3. Falta de conocimiento & 7. Falta de recursos & 11. Falta de conciencia situacional \\ \hline
4. Distracción & 8. Presión & 12. Normas (malas costumbres del grupo) \\ \hline
\end{tabular}
\end{center}

% ------------------------------------------------------------
\subsubsection*{2. Organización}
El curso se divide en grupos de 4 a 6 integrantes. Cada grupo trabaja sobre su propio
equipo y se divide en dos ``turnos'':
\begin{itemize}
  \item \textbf{Turno A (saliente):} comienza el trabajo y lo deja a medias.
  \item \textbf{Turno B (entrante):} recibe el trabajo y debe terminarlo.
  \item \textbf{Un observador} por grupo (puede rotar): no participa del trabajo, registra lo
        que ocurre en la Planilla de registro.
\end{itemize}
El docente entrega al Turno A una \textbf{orden de trabajo} escrita, que el Turno B
\textbf{no ve}. Ejemplo de orden de trabajo:
\begin{enumerate}
  \item Retirar la tapa superior del equipo.
  \item Desconectar los conectores J1 y J3.
  \item Retirar el módulo interno indicado por el docente.
  \item Volver a instalar el módulo, reconectar J1 y J3 y colocar la tapa con todos sus
        tornillos.
  \item Verificar el funcionamiento.
\end{enumerate}

% ------------------------------------------------------------
\subsubsection*{3. Ronda 1: pase de turno ``informal'' (20 minutos)}
\begin{enumerate}
  \item El Turno A realiza los pasos 1 a 3 y, antes de terminar, el docente anuncia el
        \textbf{fin del turno}. El docente puede agregar, sin avisar, alguna condición extra
        (por ejemplo, dejar un tornillo de la tapa colocado sin ajustar).
  \item El Turno A tiene \textbf{un minuto} para contarle al Turno B, \textbf{solo de palabra},
        en qué quedó el trabajo. No pueden escribir ni mostrar el equipo.
  \item El Turno B termina el trabajo, \textbf{sin ver la orden de trabajo} y sin hacer
        preguntas al Turno A.
  \item El observador registra todo en la Planilla de registro.
\end{enumerate}

% ------------------------------------------------------------
\subsubsection*{4. Ronda 2: pase de turno con planilla (25 minutos)}
Se repite el ejercicio sobre el mismo equipo (o se intercambian los equipos entre grupos), y
\textbf{se invierten los roles}: quienes fueron Turno B ahora son Turno A. Esta vez:
\begin{enumerate}
  \item El Turno A completa la \textbf{Planilla de pase de turno} antes de irse y
        \textbf{señaliza} en el equipo lo que quedó abierto o desconectado.
  \item El pase se hace \textbf{cara a cara, frente al equipo}, leyendo la planilla. El Turno B
        puede hacer preguntas y firma que recibió el trabajo.
  \item El Turno B termina la tarea usando la planilla y la orden de trabajo, y realiza una
        \textbf{verificación final} antes de cerrar.
  \item El observador vuelve a completar la Planilla de registro.
\end{enumerate}

% ------------------------------------------------------------
\subsubsection*{5. Puesta en común (20 minutos)}
Cada grupo compara los resultados de las dos rondas y responde:
\begin{enumerate}
  \item ¿Cuántos errores u omisiones hubo en cada ronda? (tornillos faltantes o flojos,
        conectores sin reconectar o cruzados, pasos omitidos, verificación no realizada).
  \item ¿Qué factores de la docena sucia aparecieron en la Ronda 1? Den un ejemplo concreto de
        cada uno.
  \item ¿Qué cambió en la Ronda 2? ¿Qué parte de la planilla fue la más útil?
  \item ¿Alguien del Turno B notó que algo no cerraba y no lo dijo? ¿Por qué? (asertividad,
        presión de tiempo).
  \item ¿Qué tienen en común sus errores con los casos de Continental Express 2574 y
        British Airways 5390?
  \item Propongan una mejora para la planilla de pase de turno del taller de la escuela.
\end{enumerate}

% ------------------------------------------------------------
\subsubsection*{6. Evaluación}
\begin{itemize}
  \item Participación y cumplimiento del rol asignado en ambas rondas.
  \item Calidad y claridad de la Planilla de pase de turno completada.
  \item Correcta identificación de los factores humanos en la puesta en común.
  \item Respuestas escritas a las preguntas del punto 5 (entrega individual o grupal).
\end{itemize}

\begin{recuadro}
\textbf{Idea central:} en mantenimiento, \textbf{lo que no está escrito no está hecho}. El
pase de turno se hace por escrito, cara a cara y frente al trabajo; ante la duda se verifica
con la documentación aprobada y nunca se supone.
\end{recuadro}

% ============================================================
% PLANILLA DE PASE DE TURNO (para imprimir)
% ============================================================
\clearpage
\begin{center}
{\large\bfseries PLANILLA DE PASE DE TURNO}\\[2pt]
{\small Taller de Aviónica -- E.E.S.T. N° 7 ``Taller Regional Quilmes''}
\end{center}

\renewcommand{\arraystretch}{1.4}
\noindent
\begin{tabular}{|>{\raggedright\arraybackslash}p{0.34\linewidth}|p{0.595\linewidth}|}
\hline
\textbf{Fecha} & \\ \hline
\textbf{Equipo / aeronave} & \\ \hline
\textbf{N° de orden de trabajo} & \\ \hline
\textbf{Turno saliente (nombres)} & \\ \hline
\textbf{Turno entrante (nombres)} & \\ \hline
\end{tabular}

\vspace{5pt}
\noindent\textbf{1. Tareas realizadas} (con el paso de la orden de trabajo)\\[2pt]
\begin{tabular}{|p{0.10\linewidth}|p{0.60\linewidth}|p{0.215\linewidth}|}
\hline
\textbf{Paso} & \textbf{Descripción de lo realizado} & \textbf{Firma} \\ \hline
 & & \\ \hline
 & & \\ \hline
 & & \\ \hline
\end{tabular}

\vspace{5pt}
\noindent\textbf{2. Tareas pendientes}\\[2pt]
\begin{tabular}{|p{0.10\linewidth}|p{0.83\linewidth}|}
\hline
\textbf{Paso} & \textbf{Qué falta hacer} \\ \hline
 & \\ \hline
 & \\ \hline
 & \\ \hline
\end{tabular}

\vspace{5pt}
\noindent\textbf{3. Estado en que queda el equipo}\\[2pt]
\begin{tabular}{|p{0.45\linewidth}|p{0.48\linewidth}|}
\hline
\textbf{Desmontado / desconectado} & \textbf{Dónde quedó (ubicación, rótulo)} \\ \hline
 & \\ \hline
 & \\ \hline
\textbf{Tornillería / fijaciones retiradas} & \textbf{Cantidad y dónde están} \\ \hline
 & \\ \hline
\end{tabular}

\vspace{5pt}
\noindent
\begin{tabular}{|p{0.62\linewidth}|c|c|}
\hline
\textbf{4. Verificaciones} & \textbf{Sí} & \textbf{No} \\ \hline
¿Se señalizó en el equipo que el trabajo está incompleto? & \hspace{0.6cm} & \hspace{0.6cm} \\ \hline
¿Hay herramientas o elementos sueltos dentro del equipo? & & \\ \hline
¿Se informó alguna anomalía o duda? (detallar abajo) & & \\ \hline
\end{tabular}

\vspace{5pt}
\noindent\textbf{5. Observaciones, anomalías o advertencias}\\[2pt]
\begin{tabular}{|p{0.955\linewidth}|}
\hline
\\ \hline
\\ \hline
\end{tabular}

\vspace{10pt}
\noindent
\begin{tabular}{|p{0.46\linewidth}|p{0.46\linewidth}|}
\hline
\textbf{Entrega (turno saliente)} & \textbf{Recibe (turno entrante)} \\
Firma y aclaración: & Firma y aclaración: \\[14pt] \hline
\end{tabular}
\renewcommand{\arraystretch}{1.15}

% ============================================================
% PLANILLA DE REGISTRO DEL OBSERVADOR (para imprimir)
% ============================================================
\clearpage
\begin{center}
{\large\bfseries PLANILLA DE REGISTRO DEL OBSERVADOR}\\[2pt]
{\small Simulación de pase de turno -- MRM}
\end{center}

\renewcommand{\arraystretch}{1.6}
\noindent
\begin{tabular}{|>{\raggedright\arraybackslash}p{0.34\linewidth}|p{0.595\linewidth}|}
\hline
\textbf{Grupo} & \\ \hline
\textbf{Observador/a} & \\ \hline
\end{tabular}

\vspace{8pt}
\noindent\textbf{Errores y omisiones detectados}\\[2pt]
\begin{tabular}{|p{0.50\linewidth}|>{\centering\arraybackslash}p{0.19\linewidth}|>{\centering\arraybackslash}p{0.19\linewidth}|}
\hline
\textbf{Tipo de error} & \textbf{Ronda 1} & \textbf{Ronda 2} \\ \hline
Tornillos faltantes, flojos o sobrantes & & \\ \hline
Conectores sin reconectar o cruzados & & \\ \hline
Pasos de la orden de trabajo omitidos & & \\ \hline
Componente mal instalado & & \\ \hline
Herramienta o elemento olvidado en el equipo & & \\ \hline
Verificación final no realizada & & \\ \hline
\textbf{Total} & & \\ \hline
\end{tabular}

\vspace{8pt}
\noindent\textbf{Factores de la docena sucia observados} (marcar y anotar un ejemplo)\\[2pt]
\begin{tabular}{|p{0.36\linewidth}|>{\centering\arraybackslash}p{0.08\linewidth}|>{\centering\arraybackslash}p{0.08\linewidth}|p{0.34\linewidth}|}
\hline
\textbf{Factor} & \textbf{R1} & \textbf{R2} & \textbf{Ejemplo observado} \\ \hline
1. Falta de comunicación & & & \\ \hline
2. Complacencia & & & \\ \hline
3. Falta de conocimiento & & & \\ \hline
4. Distracción & & & \\ \hline
5. Falta de trabajo en equipo & & & \\ \hline
6. Fatiga & & & \\ \hline
7. Falta de recursos & & & \\ \hline
8. Presión & & & \\ \hline
9. Falta de asertividad & & & \\ \hline
10. Estrés & & & \\ \hline
11. Falta de conciencia situacional & & & \\ \hline
12. Normas & & & \\ \hline
\end{tabular}

\vspace{8pt}
\noindent\textbf{Comentarios del observador}\\[2pt]
\begin{tabular}{|p{0.955\linewidth}|}
\hline
\\ \hline
\\ \hline
\\ \hline
\end{tabular}
\renewcommand{\arraystretch}{1.15}
